\section{Conclusão}

Neste trabalho, foi realizado o \textit{fine-tuning} do \ac{mllm} \ac{llama}-3.2-11B com as técnicas de treinamento
\ac{qlora} e \ac{lora}. Os modelos provenientes do \textit{fine-tuning} nas duas etapas de desenvolvimento foram
analisados e comparados com o modelo base e também entre si, oferecendo informações relevantes sobre o impacto dos
diferentes métodos de treinamento.

Os resultados apresentados constataram que o \textit{fine-tuning} do modelo \ac{llama}-3.2-11B pode ser feito de forma
eficiente através do treinamento com \ac{qlora}. Os modelos treinados com este método obtiveram um desempenho
semelhante ao dos modelos treinados com \ac{lora} e requisitaram menos memória durante o treinamento.

Na etapa inicial, os resultados obtidos foram satisfatórios e demonstraram que através do \textit{fine-tuning} com
\ac{qlora} com 5000 amostras do conjunto de dados é possível atingir uma acurácia na classificação
de lesões de pele de 85,4\%, o que é próximo da acurácia de 86,2\% obtida com o \ac{lora} para o mesmo volume de dados.
Porém, os resultados da etapa final foram poucos satisfatórios quanto ao critério de classificação de lesões, atingindo
apenas 45,2\% de acurácia após o treinamento com \ac{qlora}. Embora o volume de dados utilizado nesta etapa tenha sido
similar, com um total de 4736 pares de imagem e texto de resposta, o conteúdo em si das imagens e do texto eram
consideravelmente diferentes. As imagens de aproximação utilizadas possuem uma complexidade visual muito maior, pois
elas incluem elementos extras além da lesão de pele, como o resto do ambiente ao redor do local da lesão. Além disso,
os textos utilizados possuem não só uma seção para a classificação da lesão, mas também uma seção para a classificação
de risco, conclusão sobre a lesão e conclusão geral.

O volume de dados no conjunto de treinamento pode ter sido insuficiente para o \textit{fine-tuning}, impossibilitando
os modelos finais de realizarem uma tarefa complexa como a geração de laudos com classificações e recomendações. Além
do volume insuficiente, o desbalanceamento do conjunto de dados, caracterizado pela raridade de algumas lesões, pode
ter impactado os modelos treinados, levando a um desempenho inferior na classificação de lesões com poucas ocorrências
no conjunto de dados. Outro aspecto negativo do conjunto foi a variabilidade dentro de algumas classes de lesões. As
classes ``Lesões de pele que necessitam de avaliação presencial com dermatologista'', ``Lesões de pele que necessitam
de avaliação presencial com dermatologista pediátrico'' e ``Outras'' compõem 27,11\% do conjunto de dados de
treinamento e incluem diversas lesões não relacionadas. Este grande volume de classes incertas pode ter levado ao baixo
desempenho observado nas tarefas de classificação.

O desempenho na geração dos textos de conclusão na etapa final foi razoavelmente satisfatório, porém, alguns problemas
foram identificados, como erros ortográficos e formatações não ideais. Essas características também são provenientes de
inconsistências e erros no conjunto de dados.

Por fim, os resultados dos testes na segunda etapa sobre o conjunto de treinamento evidenciaram que o
\ac{llama}-3.2-11B após o \textit{fine-tuning} não foi capaz de capturar o padrão de classificação de lesões de pele,
o que indica que para este tipo de tarefa, um conjunto de dados limpo e bem estruturado é necessário.

Em geral, considera-se necessária a adaptação e expansão do conjunto de dados para que o desenvolvimento de um novo
modelo com um maior desempenho seja possível.

\subsection{Trabalhos Futuros}

Considerando os resultados e limitações constatadas neste projeto, propõe-se a construção de um novo conjunto de dados
com um maior volume e consistência entre as classes. Seria necessário incluir também uma nova classe referente a casos
saudáveis, pois atualmente o modelo passa pelo \textit{fine-tuning} para classificar somente diferentes lesões de pele,
desconsiderando a possibilidade da ausência de lesões.

Recomenda-se também a otimização aprofundada dos hiperparâmetros de treinamento, possivelmente com técnicas de busca em
grade e ferramentas como W\&B Sweeps e Optuna \cite{wab_sweeps, optuna_2019}.

A utilização de outros \acp{mllm} também é recomendada, particularmente a do Gemma 3 de \textcite{team2025gemma}. Isso
tornaria possível a realização de uma comparação mais avançada entre os diferentes modelos.

Por fim, sugere-se a utilização de avaliações manuais sobre partes dos resultados dos testes. Atualmente, os resultados
são analisados de forma automática, o que pode levar à desconsideração de aspectos importantes do texto gerado. Com a
participação de dermatologistas qualificados, seria possível avaliar o conteúdo destes textos de forma mais
aprofundada.
